% status: FULL
% Printed Chapter 32. Stable source/lab ID: ch15-continuous-batching.
% Preserves measurements from research/measurements/2026-09-24-m1-part3.md
% and research/measurements/2026-09-23-m1-llama32-3b.md.
\chapter{Continuous Batching}\label{ch:continuous-batching}

\begin{ChapterIntent}
A request lives for many model iterations, but a batch should not. Continuous
batching rebuilds the executable work set at iteration boundaries so that
finished requests leave, newly admissible requests enter, and every scheduled
token is charged against explicit sequence, token, and memory limits. This
chapter follows four requests through fixed request batching and
iteration-level batching, derives the scheduler state needed to keep ragged
requests correct, and builds a deterministic CPU scheduler whose accounting can
be checked independently. The goal is not a particular engine policy. It is
the set of invariants that any production policy must preserve.
\end{ChapterIntent}

\OpeningSection{The request that waited for a stranger's essay}

Start llama.cpp's server on the recorded Apple M1 setup with four request
slots and Llama~3.2~3B. One client asks for a 512-token answer. Two seconds
later a second client asks for 32 tokens. With continuous batching disabled,
the second client's first token arrived \textbf{20.5 to 21.6 seconds} after
submission in three runs. With continuous batching enabled, its first token
arrived in \textbf{0.50 seconds}, and its complete answer arrived 2.9 seconds
after submission (measured; record
\texttt{research/measurements/2026-09-24-m1-part3.md}).

That difference is not a faster matrix multiplication. It is a different
answer to a scheduling question: \emph{when may the membership of the next
forward pass change?} Fixed request batching says ``after the current cohort
finishes.'' Continuous batching says ``at an iteration boundary, subject to
capacity and ownership rules.'' The second answer turns a batch from a durable
container into a temporary execution plan.

The first request did not receive the improvement for free. In the same
recorded experiment, its long answer took 23.3--25.7 seconds with the newcomer
instead of 22.1--23.2 seconds without it. During the shared interval, a step
took about 75~ms rather than about 40~ms. On that M1 workload, a two-sequence
step was therefore close to the cost of two one-sequence steps. Continuous
batching bought latency isolation for the newcomer, not a large aggregate
throughput gain. On hardware where a larger decode batch reuses weights much
more efficiently, the same scheduling mechanism can buy both latency and
throughput. The mechanism does not guarantee either outcome by itself.

This chapter starts where Chapter~\ref{ch:request-lifecycle} stopped. A request
already has a stable identity, a terminal-state owner, a cancellation path and
a retirement rule. The scheduler must not weaken those guarantees merely
because it wants to reuse a row quickly.

\Foundations{A batch row is not a request}{
A request is durable control state. A batch row is a temporary assignment for
one executable step. The row may contain prompt tokens for one request, one
decode token for another, and disappear completely on the next iteration.
Request identity, output ownership and cache ownership therefore cannot be
encoded only by the row index. Chapter~\ref{ch:prefill-decode} introduced the
different arithmetic shapes of prefill and decode; Chapter~\ref{ch:request-lifecycle}
established that terminal publication and physical retirement are separate
events. We now combine those ideas in a scheduler.
}

\NapkinSection{Napkin math: why request-level batches waste slots}

Suppose a fixed cohort contains $B$ requests and request $i$ needs $L_i$
decode iterations. If the implementation keeps all $B$ rows until the longest
request finishes, it reserves
\[
  W_{\mathrm{reserved}} = B\max_i L_i
\]
sequence-steps while useful decode work is only
\[
  W_{\mathrm{useful}} = \sum_{i=1}^{B} L_i .
\]
The sequence-slot utilization is therefore
\begin{equation}
  \eta_{\mathrm{static}}
  = \frac{\sum_i L_i}{B\max_i L_i}
  = \frac{\bar L}{L_{\max}} .
  \label{eq:static-occupancy}
\end{equation}
These counts are exact under the stated convention: one active request for one
decode iteration is one sequence-step. They are not FLOP counts and they do
not predict wall-clock speed without a step-time model.

Take four illustrative output budgets:
\[
  (L_A,L_B,L_C,L_D)=(6,2,4,1).
\]
A fixed four-request cohort occupies $4\times6=24$ sequence-steps but performs
$6+2+4+1=13$ useful sequence-steps, so
\[
  \eta_{\mathrm{static}}=13/24\approx0.542.
\]
Almost half of the reserved row-time is empty after shorter requests finish.
The deeper problem is temporal: a fifth request arriving while $A$ is still
running cannot use one of those empty logical rows if cohort membership is
frozen.

Continuous batching removes that particular restriction. It does not create
unlimited capacity. If the scheduler allows at most $S_{\max}$ active
sequences and at most $T_{\max}$ scheduled tokens in an iteration, a candidate
plan must satisfy
\begin{align}
  |\mathcal A_k| &\le S_{\max},\\
  \sum_{r\in\mathcal A_k} q_{r,k} &\le T_{\max},
  \label{eq:budgets}
\end{align}
where $\mathcal A_k$ is the set of requests scheduled at iteration $k$ and
$q_{r,k}$ is the number of input positions request $r$ contributes to that
iteration. For ordinary decode, $q_{r,k}=1$. For prefill, $q_{r,k}$ can be
many prompt positions. Equation~\eqref{eq:budgets} is why ``batch size'' is
ambiguous in modern inference: a sequence budget and a token budget constrain
different resources.

\section{One cohort, two schedulers}

We will follow four illustrative requests. Time is measured in scheduler
ticks, not milliseconds, so the example exposes policy rather than pretending
to be a hardware benchmark.

\begin{center}
\begin{tabular}{lrrrr}
\toprule
Request & arrival & prompt tokens & output budget & cancel tick\\
\midrule
A & 0 & 3 & 4 & --\\
B & 0 & 1 & 2 & --\\
C & 1 & 2 & 3 & --\\
D & 3 & 1 & 2 & 5\\
\bottomrule
\end{tabular}
\end{center}

Assume $S_{\max}=2$ active requests and $T_{\max}=3$ scheduled tokens per
iteration. To keep this chapter focused, a request's prefill may be scheduled
as one chunk only when its remaining prompt fits the current token budget.
Chapter~\ref{ch:chunked-prefill} will remove that simplifying restriction.
Likewise, the scheduler asks a cache allocator whether a request fits, but
Chapter~\ref{ch:paged-kv} owns the block-allocation mechanism.

With \emph{fixed request batching}, $A$ and $B$ form the first cohort because
both arrive at tick 0. Their membership is frozen. Even after $B$ finishes,
$C$ cannot enter until $A$ also finishes. The empty place beside $A$ is wasted.
Only then may $C$ and $D$ form another cohort. If $D$ is cancelled while that
second cohort runs, its terminal result may be published, but its physical
slot cannot be reused until outstanding work carrying its lease has retired.

With \emph{iteration-level batching}, membership is reconsidered after every
completed step. When $B$ retires, $C$ can become active while $A$ continues.
The important word is \emph{retire}, not merely \emph{finish}. Chapter 31's
ownership rule survives intact: the scheduler may observe a terminal decision
before every device or worker result associated with the old lease is gone.

\EngineFigure{ch15-moving-cohort}{A separate three-request decode-only
illustration, not the four-request lab trace. A occupies six ticks; B occupies
two. Fixed request batching leaves a hole after B finishes;
iteration-level batching admits C at the next safe boundary. The figure models
scheduler ticks, not measured wall-clock time.}{fig:ch15-moving-cohort}

\section{Admission is not scheduling}

Two decisions that are often conflated must be kept separate.

\emph{Admission} decides whether a request may enter the engine's managed
working set. It is a capacity and service-protection decision. An engine may
reject or queue a request because too many requests are already in flight,
because its prompt would exceed a configured backlog limit, because the KV
allocator cannot reserve safe capacity, or because a tenant has exhausted a
quota. Admission can happen before the request has ever appeared in a model
batch.

\emph{Scheduling} decides which admitted requests contribute work to the
\emph{next} iteration. A request can be admitted but waiting. It can be active
but unscheduled for one iteration. It can be terminal but not retired. A
scheduler that stores only ``in batch / not in batch'' has too few states to
express these cases safely.

For the teaching implementation we use
\[
 \textsc{Queued}\rightarrow\textsc{Admitted}\rightarrow
 \textsc{Running}\rightarrow\textsc{Terminal}\rightarrow\textsc{Retired}.
\]
Cancellation can move a queued, admitted or running request to
\textsc{Terminal}. Only retirement makes its reusable slot available. A
production engine may split these states further, but it must preserve the
same ownership distinction.

\EngineFigure{ch15-request-state}{Conceptual state machine behind the teaching
scheduler. Terminal publication is irreversible; physical capacity returns
only on retirement. A generation number guards reused slots against late
results.}{fig:ch15-request-state}

This separation gives backpressure a precise location. If the waiting queue is
unbounded, continuous batching can keep the accelerator busy while user
latency diverges. A saturated engine that accepts every request is not
healthy; it is accumulating future latency. Capacity limits must therefore
bound some combination of queued requests, queued prompt tokens, active
sequences and cache residency. A rejection such as HTTP 503 can be a more
honest service response than accepting work whose latency objective is already
impossible.

\section{Prefill and decode consume different budgets}

Autoregressive inference has two recurring shapes. During \emph{prefill}, a
request contributes many known prompt positions. During \emph{decode}, it
usually contributes one new position per iteration. Both use the same model
weights, but their arithmetic intensity, attention work and latency role are
different.

Consider $A$ with three prompt tokens and $B$ with one. With
$T_{\max}=3$, scheduling all of $A$'s prompt consumes the entire token budget.
Although the sequence budget still has room for $B$, the token budget does
not. On the next iteration $A$ may need one decode token while $B$ contributes
one prefill token, for only two scheduled tokens total. A sequence-only limit
cannot express this distinction.

This also explains why ``continuous batching'' is not synonymous with
``decode batching.'' The active set may be ragged in both phase and length.
A production step can contain decode work for several requests and prompt work
for another if its kernels and scheduler support that combination. The
scheduler needs an interface that says how many positions each request may
advance and a cache interface that says whether those positions can be made
resident. Detailed chunk selection belongs to
Chapter~\ref{ch:chunked-prefill}; detailed block placement belongs to
Chapter~\ref{ch:paged-kv}.

A useful representation is a per-request work ticket
\[
  w_{r,k}=(\mathrm{id}_r,g_r,p_r,c_r,o_r,q_{r,k}),
\]
where $\mathrm{id}_r$ is stable request identity, $g_r$ is the slot generation,
$p_r$ is prompt length, $c_r$ is the number of model positions already
computed, $o_r$ is the number of committed output tokens, and $q_{r,k}$ is
the number of positions authorized for this iteration. The batch builder may
pack these tickets into any row order. Correctness must not depend on that
order.

\section{Ragged batches need explicit metadata}

Suppose the next executable step contains
\[
 q_A=1,\qquad q_B=1,\qquad q_C=2.
\]
There are four token rows but three requests. A flattened token tensor can be
written as
\[
  \mathbf X\in\mathbb R^{4\times d},
\]
where $d$ is model width. The rows might be ordered
\[
  [A_3,\ B_1,\ C_0,\ C_1].
\]
The subscript is the logical position within that request. The model-wide
linear layers can treat these as four ordinary rows. Attention cannot. It
must know that $A_3$ may attend to $A$'s earlier positions, $B_1$ to $B$'s,
and the two $C$ rows to the correct causal prefix of $C$.

One compact host representation uses request offsets
\[
  \mathbf o=[0,1,2,4].
\]
Request $j$ owns flattened rows
$[\mathbf o_j,\mathbf o_{j+1})$. The differences
$(1,1,2)$ recover $q_A,q_B,q_C$. A second vector carries logical positions,
for example
\[
  \mathbf p=[3,1,0,1].
\]
A cache mapping associates each request and logical position with storage.
The next chapter will make that mapping paged; this chapter requires only the
interface.

At four-byte integers, the illustrative offsets occupy $4\times4=16$ bytes
and positions another $4\times4=16$ bytes. Those byte counts are derived, not
measured. Real runners carry more metadata: cache block tables, sequence
lengths, sampling state, adapter identifiers, grammar state, multimodal
metadata and often device-resident persistent buffers. The important
engineering principle is that raggedness is represented explicitly rather
than hidden in padding or inferred from row number.

\EngineFigure{ch15-ragged-batch}{Memory view of an illustrative ragged batch.
Four flattened token rows belong to three durable requests; offsets and
logical positions recover ownership while row numbers remain temporary. Cache
addresses are intentionally abstract because Chapter 33 owns paging.}
{fig:ch15-ragged-batch}

\section{The iteration transaction}

A safe scheduler iteration is easiest to reason about as a transaction with
five boundaries.

First, \emph{collect events}. Arrivals, cancellation observations and worker
completions are applied to durable request records. A terminal decision is
monotonic: once cancellation or completion owns the outcome, a later event
cannot replace it.

Second, \emph{retire safe work}. A terminal request returns physical capacity
only when every outstanding work ticket for its current generation has
retired. This is the Chapter 31 invariant that prevents a late result from
being published into a reused request.

Third, \emph{admit}. The engine moves queued requests into its managed active
set only while request, token-backlog and memory limits permit. Admission
policy can be FCFS, priority-aware or tenant-aware; continuous batching does
not choose among them.

Fourth, \emph{schedule}. The scheduler walks eligible requests and grants
$q_{r,k}$ positions while both sequence and token budgets remain. The teaching
lab uses FCFS admission and deterministic physical-slot order for scheduling.
A production scheduler may reserve decode
capacity, cap long prefills, preempt work or use deadlines. Those policies
change latency and fairness but not the accounting invariants.

Fifth, \emph{launch}. The batch builder snapshots request identity and
generation into immutable work tickets. The worker result must return the same
pair. When the result comes back, the owner validates
$(\mathrm{id},g)$ before committing it.

\EngineFigure{ch15-iteration-sequence}{Sequence diagram for one scheduler
iteration. Events and retirement precede admission; admission precedes
scheduling; launch snapshots identity plus generation; result publication
validates the snapshot before mutating request state.}{fig:ch15-iteration-sequence}

The transaction need not be implemented under one mutex. CPU scheduling can
overlap accelerator execution. The logical dependency, however, must remain:
a plan cannot consume capacity that has not safely returned, and a result
cannot mutate state whose generation no longer matches.

\section{Safe completion, cancellation and slot reuse}

Assume request $D$ occupies physical slot 1 with generation 7. The scheduler
launches work ticket $(D,7)$, then observes cancellation. The request becomes
terminal immediately from the client's perspective. The engine may stop
scheduling new work for it, but ticket $(D,7)$ is still outstanding.

The tempting optimization is to assign slot 1 to request $E$ immediately.
Suppose the slot's request pointer now names $E$. When the old result returns,
code that validates only the slot index sees a valid slot and may append
$D$'s token to $E$. This is an ABA-style reuse bug: the location looks valid
again even though its ownership epoch changed.

The fix is a generation or lease identity. Reuse increments the generation,
so $E$ owns $(1,8)$. A late $(D,7)$ result cannot match it. The stronger rule
used in the teaching lab is to avoid reuse until outstanding work reaches
zero; generation checking remains defense in depth and makes the failure
observable.

Three invariants are worth writing as assertions:

\begin{enumerate}
\item \textbf{Budget invariant:}
$\sum_r q_{r,k}\le T_{\max}$ and the number of requests with
$q_{r,k}>0$ is at most $S_{\max}$.
\item \textbf{Terminal monotonicity:} once a request is terminal, no later
worker result changes its externally visible terminal reason.
\item \textbf{Lease validity:} a worker result may mutate request state only
when its request identity and generation match an outstanding ticket.
\end{enumerate}

These are correctness properties. They should be tested independently of
performance.

\section{Capacity and backpressure}

Continuous batching improves utilization only while there is schedulable work
that fits. It cannot manufacture KV memory, host queue capacity or latency
budget.

Let $Q$ be the number of admitted-but-waiting requests and $\lambda$ the
arrival rate in requests/s. If the engine's sustainable completion rate is
$\mu<\lambda$, then $Q$ grows without bound in an unbounded queue. The
accelerator may show excellent utilization while TTFT becomes unacceptable.
That is why admission control belongs before scheduling.

A token backlog can be more informative than a request count. Ten requests
with 32-token prompts and ten requests with 32,000-token prompts are not the
same prefill obligation. If the engine estimates sustainable prefill service
at $R_p$ prompt tokens/s and wants a prefill-queue contribution to TTFT below
$\tau$, a rough derived bound is
\[
  Q_{\mathrm{prompt}}\lesssim R_p\tau
\]
prompt tokens. This is a capacity-planning relation, not a universal
admission formula. Cache hits, chunking, multimodal encoders and mixed decode
work change the realized service rate.

Backpressure is therefore part of the continuous-batching story even though
it is not part of the batching algorithm. Without it, the scheduler can
optimize device occupancy while the service violates its user-facing target.

\section{Throughput, TTFT, inter-token latency and goodput}

Four metrics answer different questions.

\emph{Throughput} counts useful tokens completed per unit time across the
engine. A larger active set often improves it while the model is
memory-bandwidth limited because more requests share a weight read.

\emph{Time to first token} (TTFT) measures submission to first generated
token. Continuous admission can reduce TTFT dramatically compared with a
frozen cohort, as the opening measurement shows. Long prefills or a deep
waiting queue can still make TTFT poor.

\emph{Inter-token latency} (ITL), or time per output token (TPOT) under a
specified convention, measures the cadence after generation starts. Larger
decode batches can increase each request's step time even while aggregate
throughput improves.

\emph{Goodput} counts only work that satisfies a service objective. If a
system produces 1,000 tokens/s but half the requests miss the TTFT or ITL
target, raw throughput hides the service failure. Continuous batching is a
mechanism for exposing more scheduling choices; goodput is one way to judge
whether those choices serve users.

The recorded M1 decode staircase illustrates the tension. Aggregate decode
throughput rose from roughly 20 tokens/s at one sequence to 176 tokens/s at
32 sequences, while per-sequence step time rose from roughly 50~ms to 182~ms
(record \texttt{research/measurements/2026-09-23-m1-llama32-3b.md}).
Those are measurements for that recorded setup, not a universal batch curve.
The scheduler needs a policy target because maximizing aggregate throughput
and minimizing per-request latency are different optimization problems.

\section{Continuous batching does not imply fairness}

Consider two tenants. Tenant X keeps 100 short decode requests ready. Tenant Y
submits one long prompt. A scheduler that always fills available token budget
with already-running decode work can keep X moving indefinitely while Y waits.
Reverse the priority and a large prefill can inflate ITL for every decoder.
Both schedulers use continuous batching. Neither policy follows from the
mechanism.

Fairness requires an additional rule: FCFS, weighted fair sharing, priority
classes, deadlines, age boosts, reserved decode budget, prefill caps or another
explicit policy. Latency objectives likewise require control. Continuous
batching says when membership \emph{may} change. It does not say whose work
\emph{should} run.

This distinction matters when reading production source. A token budget is a
mechanism. A long-prefill threshold is policy. A queue cap is admission. A
cache watermark is capacity protection. They interact, but collapsing them
into one word---``batching''---makes failures difficult to diagnose.

\section{What the teaching implementation builds}

The lab in \texttt{code/labs/ch15-continuous-batching/} is a deterministic
CPU state machine. It does not run a Transformer. That is intentional: the
scheduler can be tested without GPU timing noise.

Each request has an immutable ID, arrival tick, prompt length and output
budget; mutable state tracks computed prompt positions, committed output,
terminal reason, physical slot, generation and outstanding tickets. Each
iteration performs the transaction described above and emits a trace:
arrivals, admissions, scheduling grants, terminal transitions, retirement and
rejected replies. Completion is a synchronous method call in this lab, not
a separate trace event; a production trace needs submission and completion
timestamps as well.

An independent oracle recomputes budget accounting from the emitted trace
rather than calling the scheduler's accounting methods. Tests cover the hand
cohort, capacity bounds, cancellation before admission, cancellation during
outstanding work, duplicate replies and late replies after slot reuse.
Input validation also rejects duplicate IDs and prompts that cannot fit
an unchunked iteration. These are contract tests, not throughput measurements.

The lab's unit of time is one tick. It demonstrates state and accounting, not
GPU performance, kernel efficiency, paged allocation, prefix reuse or
chunked-prefill quality. Those boundaries are part of the lesson.

\section{Walk the trace before changing the policy}

The most useful way to read a scheduler is to explain each change of
ownership. Start with the trace from the executable, not a proposed
throughput improvement. At tick zero both A and B are admitted: admission
consumes two request slots. Only A is scheduled: its three prompt positions
consume the entire token budget. B therefore owns a slot without receiving
any model work yet. This distinction is not an implementation accident.
It is the difference between reserving the right to run and actually
running. Monitoring only the number of admitted requests would miss it.

At tick one C arrives and waits. A requires one decode position; B still
requires its complete one-position prompt. Together they consume two of
the three available positions. The remaining position is not an invitation
to admit C: both request slots are already owned. At tick two A and B each
receive one decode position. At tick three D arrives while B receives its
second and final output position. B becomes terminal at the end of that
iteration. Its terminal event is a statement about the request's outcome,
not an admission grant for another request.

At tick four retirement releases B's slot. C, which arrived before D,
inherits that physical location with a different generation. A receives
one decode position and finishes; C receives its two-position prompt.
The total is exactly three positions, but the rows have different roles.
The model-side batch builder must know which row reads an existing cache,
which rows extend a prompt, and which sequence owns every position.
A single rectangular hidden-state matrix does not erase those distinctions.

At tick five D is cancelled while it is still waiting. It never owned a
slot and never launched a ticket. There is consequently no slot-retirement
event for D. A retires, but the cancelled waiting entry is skipped rather
than admitted. C decodes at ticks five, six and seven, then retires at tick
eight. This small trace exposes a tempting test mistake: requiring every
cancelled request to retire a slot would invent ownership it never had.
The test for active cancellation must create a genuinely outstanding
ticket, observe cancellation, and prove that retirement waits for that
ticket's completion.

\begin{center}
\begin{tabular}{@{}clll@{}}
\toprule
Tick & Granted work & Total positions & Boundary consequence\\
\midrule
0 & A prompt: 3 & 3 & B admitted but not scheduled\\
1 & A decode: 1; B prompt: 1 & 2 & C waits for a slot\\
2--3 & A decode: 1; B decode: 1 & 2 each & B finishes at tick 3\\
4 & A decode: 1; C prompt: 2 & 3 & B retired; A finishes\\
5 & C decode: 1 & 1 & D cancelled without admission\\
6--7 & C decode: 1 & 1 each & C finishes at tick 7\\
8 & None & 0 & C's slot retires\\
\bottomrule
\end{tabular}
\end{center}

These are deterministic logical ticks. They are not milliseconds, and
three positions do not necessarily take three times as long as one.
The lab counts a prompt pass separately from output-token commits; it
does not implement the usual optimization in which the final prompt row's
logits select the first output token. When connecting it to a real decoder,
define that boundary explicitly before comparing the number of model calls.
Otherwise a supposedly faster schedule can simply be counting a different
amount of work.

\section{A reply needs more than the right slot}

Imagine an output message addressed to apartment four. The address tells
the delivery system where to go, but not whether the resident is still
the intended recipient. A generation distinguishes successive occupants
of a physical slot. A ticket serial distinguishes successive pieces of
work for the same occupant. Both checks matter. If a worker retries a
reply before the request retires, the request ID and generation can still
match. Accepting that duplicate could decrement an outstanding counter
twice or append the same output twice.

The teaching implementation stores each pending ticket by serial.
A reply must match the current slot owner and the complete pending
ticket, including its token grant. The first valid reply consumes the
pending entry. A second copy is rejected without changing the request.
The pending map is not a general network exactly-once protocol: it is
an in-memory ownership check within one scheduler instance. Recovering
after a process crash would additionally require durable request epochs,
replay policy and a definition of what the client has already observed.

Cancellation changes whether a valid reply may publish output, but it
does not pretend that the reply never existed. A matching completion for
a cancelled request retires outstanding ownership while discarding its
payload. Only then can the slot be released. In a GPU implementation
this boundary must correspond to actual completion of all consumers of
the memory, often established by an event or a dependency on a stream.
A generation check on the CPU cannot repair a device write that already
overwrote another request's KV bytes.

There is also a budget-replay hazard. Calling a planning method twice
with a fresh local budget in the same logical iteration can spend the
iteration's allowance twice. This lab therefore permits one plan per
tick and skips requests with outstanding work. Overlapped production
engines may allow several launches, but then their shared budget and
computed-position frontiers must include all in-flight grants. The right
extension is an explicit reservation model, not removal of the guard.

\section{Read the batch as a data contract}

The scheduler produces a plan, not a tensor. The batch builder translates
that plan into token IDs, absolute positions, sequence offsets, cache
locations and row-to-request ownership. For a ragged batch with grants
$(1,1,2)$, the offset vector is $(0,1,2,4)$. Row two belongs to the third
request even though its physical request slot need not be two. Row three
belongs to that same request, at the next logical position. Sampling
must consume the appropriate final row's logits rather than every row
indiscriminately.

Keep request identity out of accidental tensor-index conventions.
An integer token ID identifies a vocabulary entry. An integer row
identifies a location in this launch's tensor. A logical position
identifies a point in a request's history. A physical cache address
identifies storage. All four can happen to contain the number seven,
yet substituting one for another is a correctness bug. Named metadata
fields and checked conversion boundaries make those bugs visible before
they become wrong text.

The same distinction applies at the binary boundary. Rust's native
struct layout is not a portable worker-message format. A distributed
worker protocol must specify field widths, byte order, version, bounds
and ticket identity rather than transmitting the memory image of a
struct. Within one process, an index into a live pending-ticket table
can be enough; across processes, that index needs a lifetime and restart
contract. No tensor optimization should silently choose that contract.

Finally, separate host and device costs when evaluating an improvement.
A larger cohort can amortize weight traffic while increasing metadata
packing, allocation, grammar checks and response dispatch. Time each
stage, then measure the overlapped critical path; adding independent
stage timings is not automatically the observed wall time. Report the
arrival process and output lengths alongside the batch histogram.
An excellent kernel benchmark can coexist with poor request latency
if admission waits, backpressure or a long prompt dominates the service.

\EngineFigure{sys-iteration-scheduler}{Activity view of a conceptual continuous-batching scheduler. Stable request IDs outlive temporary batch rows. Completion must establish that consumers are finished before their state is reclaimed; overlapped implementations preserve that dependency.}{fig:sys-iteration-scheduler}

\LedgerSection

\begin{ConstraintLedger}
Memory bandwidth & buys & More sequences can share each weight read when the hardware benefits from a larger token matrix. \\
Memory capacity & spends & Every admitted/running request retains request state and, in a real engine, KV-cache capacity; batching cannot exceed residency. \\
Compute & moves & Filling empty decode slots raises useful work, but mixed prefill/decode work changes the arithmetic shape and can move the bottleneck. \\
Latency & buys & New requests can enter at an iteration boundary instead of waiting for a frozen cohort to drain. \\
Latency & spends & Larger active sets can lengthen each model step; queueing and long prefills can still violate TTFT or ITL targets. \\
Correctness & risks & Row reuse, cancellation races and stale worker results can cross request ownership unless identity and generation are validated. \\
Cost & buys & Higher utilization can improve useful tokens per accelerator-dollar only when the added work still meets the service objective. \\
\end{ConstraintLedger}

\LabSection{a deterministic scheduler with an independent accounting oracle}

The lab is self-contained Rust with no external crates. Before running it,
predict the hand cohort's first two scheduling decisions. Then run:

\begin{verbatim}
cd code/labs/ch15-continuous-batching
cargo test
cargo run -- trace
cargo run -- broken late-result
\end{verbatim}

The normal trace must preserve both budgets, terminal monotonicity and lease
validity. The broken late-result mode is a non-zero instruction stub, not an
executed mutation experiment: use the focused stale-reply test, then change
its owner validation locally to observe the failure. Do not report the stub
as evidence that a broken implementation was run. The README
contains the complete predict/run/explain/break route.

\HappenedSection{from iteration-level scheduling to token-budget schedulers}

Orca named iteration-level scheduling and selective batching in OSDI 2022
\cite{yu2022orca}. Its key observation remains durable: autoregressive
requests have multiple iterations, so freezing a request cohort wastes
finished capacity and delays new arrivals. The paper reports a 36.9$\times$
throughput improvement over the compared FasterTransformer configuration at
the same latency for its GPT-3 175B evaluation. That is the authors'
measurement, not a result reproduced by this book.

The modern scheduler interface is broader. In vLLM V1 source inspected on
2026-10-05, the scheduler does not need a hard semantic split called
``prefill phase'' versus ``decode phase'' for its core token accounting.
A request carries how many tokens have been computed and how many need to be
computed; each iteration grants tokens so the first count can catch the
second. The scheduler separately maintains a per-iteration token budget.
The configuration also distinguishes maximum batched tokens, maximum
sequences, and an optional lower maximum number of active sequences. These are
source observations from the inspected revision, not claims about every vLLM
release or deployment.

That representation is important because the industrial scheduler is no
longer merely replacing finished decode rows. The same accounting boundary
must coexist with prompt chunks, speculative tokens, cache allocation,
multimodal encoder work and asynchronous execution. The teaching lab stops
before those mechanisms, but its work ticket---stable request identity,
generation and granted token count---is deliberately shaped to survive that
expansion.

Hermon's historical evidence makes a complementary point. Its earlier
concurrent runtime used several independent contexts, so concurrent requests
performed separate forward passes and separately read model weights. In the
record already preserved by this book, four concurrent 128-token requests on
an Apple M3 Pro produced 19.1 tokens/s in aggregate, while a replacement that
placed several sequences into one shared forward step produced 41.4 tokens/s.
Those are the developers' recorded measurements from the cited historical
Hermon material, not measurements rerun for this chapter. The mechanism-level
lesson is narrower than the number: concurrency is not batching unless the
requests share executable model work.

\FrontierSection{2026-10-05}

The frontier has moved from ``can requests join between iterations?'' to
``how much heterogeneous work should each iteration admit without destroying
latency or capacity?'' The vLLM scheduler source checked on 2026-10-05 exposes
distinct token, sequence and active-request controls and includes policy
controls for long prefills and queue/backlog protection. Its core scheduler
represents work as per-request scheduled token counts, allowing the same
accounting structure to cover several inference mechanisms. These are shipped
source interfaces, not evidence that every option is enabled by default;
deployment behavior still depends on configuration.

SGLang-family scheduler documentation and source inspected the same day also
describe waiting/running queues, dynamic batch construction, KV allocation and
overlap of CPU scheduling with device execution. The important trend is not a
single winning policy. It is the increasing separation of durable request
state, capacity management, scheduling policy and device execution so that
each can evolve without making batch-row identity a correctness primitive.

The next chapters fill in two interfaces this chapter deliberately leaves
abstract. Chapter~\ref{ch:paged-kv} explains how logical request positions
obtain physical KV storage without requiring one contiguous maximum-size
allocation. The later chunked-prefill chapter explains how a long prompt can
consume only part of an iteration's token budget instead of monopolizing it.

\ExercisesSection

\begin{enumerate}
\item \emph{Prediction.} With $S_{\max}=2$ and $T_{\max}=3$, requests
$A=(p=3,o=4)$ and $B=(p=1,o=2)$ arrive together. Under the lab's rule that a
prefill is not split, which request can consume the first iteration? What
budget remains?

\item \emph{Derivation.} A fixed batch has output lengths
$(2,2,2,10)$. Compute $\eta_{\mathrm{static}}$. Then compute the number of
unused sequence-steps.

\item \emph{Representation.} A ragged step grants token counts $(1,3,2)$ to
requests $(R_0,R_1,R_2)$. Construct the offset vector and state the shape of
the flattened hidden-state matrix for model width $d=4096$.

\item \emph{Correctness.} Slot 4 generation 11 belongs to request $X$ when a
worker ticket launches. $X$ is cancelled; a broken scheduler reuses slot 4
for $Y$ at generation 12. State the minimum fields a returning result must
carry to prevent publication into $Y$.

\item \emph{Implementation.} Extend the lab with a reserved decode-token
budget. Define the invariant that prevents prefills from consuming the
reservation when at least one decoder is eligible. Add a deterministic test
where the reservation changes TTFT/ITL trade-offs.

\item \emph{Falsification.} Construct a workload where continuous batching
preserves high accelerator occupancy but user goodput is zero under a stated
TTFT target. Explain why the example falsifies the claim ``high utilization
implies a healthy service.''
\end{enumerate}

\Needspace{18\baselineskip}
\section{Worked problems}

\Needspace{12\baselineskip}
\subsection{Prediction: the first iteration}
\textbf{Problem.} \emph{Prediction.} With $S_{\max}=2$ and $T_{\max}=3$, requests
$A=(p=3,o=4)$ and $B=(p=1,o=2)$ arrive together. Under the lab's rule that a
prefill is not split, which request can consume the first iteration? What
budget remains?

\textbf{Solution.} $A$'s unsplit prefill needs all three available token positions, so the
deterministic FCFS teaching scheduler grants $q_A=3$ and $q_B=0$ in the first
iteration. The sequence budget has one nominal place left, but the token budget
has zero tokens left. This is the smallest example showing why a sequence cap
cannot replace a token cap.

\Needspace{12\baselineskip}
\subsection{Derivation: fixed-batch waste}
\textbf{Problem.} \emph{Derivation.} A fixed batch has output lengths
$(2,2,2,10)$. Compute $\eta_{\mathrm{static}}$. Then compute the number of
unused sequence-steps.

\textbf{Solution.} For $(2,2,2,10)$, useful sequence-steps are $16$ and reserved sequence-steps
are $4\times10=40$. Thus $\eta_{\mathrm{static}}=16/40=0.4$. The unused
reservation is $40-16=24$ sequence-steps.

\Needspace{12\baselineskip}
\subsection{Representation: ragged offsets}
\textbf{Problem.} \emph{Representation.} A ragged step grants token counts $(1,3,2)$ to
requests $(R_0,R_1,R_2)$. Construct the offset vector and state the shape of
the flattened hidden-state matrix for model width $d=4096$.

\textbf{Solution.} Counts $(1,3,2)$ have prefix sums $(0,1,4,6)$, so the offset vector is
$\mathbf o=[0,1,4,6]$. There are six flattened token rows. With model width
$d=4096$, the hidden-state matrix is
$\mathbf X\in\mathbb R^{6\times4096}$.

\Needspace{12\baselineskip}
\subsection{Correctness: stale result}
\textbf{Problem.} \emph{Correctness.} Slot 4 generation 11 belongs to request $X$ when a
worker ticket launches. $X$ is cancelled; a broken scheduler reuses slot 4
for $Y$ at generation 12. State the minimum fields a returning result must
carry to prevent publication into $Y$.

\textbf{Solution.} The returning result must identify the durable request and its ownership
generation, here $(X,11)$, rather than only physical slot 4. The current slot
owner is $(Y,12)$, so the result cannot mutate $Y$. A system that also tracks
a unique work-ticket ID can reject duplicate completion as well.

\Needspace{12\baselineskip}
\subsection{Falsification: utilization is not goodput}
\textbf{Problem.} \emph{Falsification.} Construct a workload where continuous batching
preserves high accelerator occupancy but user goodput is zero under a stated
TTFT target. Explain why the example falsifies the claim ``high utilization
implies a healthy service.''

\textbf{Solution.} Let the TTFT objective be 100 ticks. Feed work faster than the engine can
complete it while keeping the runnable set non-empty. Accelerator utilization
can remain 100\%, yet after the queue exceeds 100 ticks every newly completed
request violates the TTFT objective. Raw utilization is therefore compatible
with zero goodput under that objective.

